\section{Рабочий проект}
\subsection{Спецификация компонентов и классов модулей веб-платформы}

\subsubsection{Спецификация класса MailMessageController}

Описание: Управляет созданием, отображением, удалением и отправкой электронных писем. Обеспечивает интеграцию с адресной книгой, управление вложениями и взаимодействие с внешним почтовым сервером через протокол JMAP. Отвечает за валидацию входящих и исходящих писем, обработку ошибок при отправке и получении, а также за маршрутизацию писем между пользователями системы.

Зависимости:
\begin{itemize}
    \item письма связаны с пользователями (отправитель и получатели) через таблицу employee;
    \item каждое письмо может содержать вложения (файлы) через связь с таблицей file;
    \item письма могут быть перемещены между пользовательскими папками через связь с таблицей folder;
    \item интеграция с ContactService для автозаполнения адресов и FileStorageService для управления вложениями.
\end{itemize}

\begin{xltabular}{\textwidth}{|p{6cm}|X|}
\caption{Данные класса MailMessageController\label{mailmessagecontroller:data}} \\ \hline
\centrow Данные & \centrow Описание \\ \hline
\endfirsthead
\continuecaption{Продолжение таблицы \ref{mailmessagecontroller:data}} 
\centrow Данные & \centrow Описание \\ \hline
\endhead
Message & Модель, представляющая письмо \\ \hline
Sender & Email отправителя (связь с таблицей employee) \\ \hline
Recipients & Список адресатов (JSON-массив) \\ \hline
Subject & Тема письма (максимум 200 символов) \\ \hline
Body & Текст письма в формате HTML/plaintext \\ \hline
Attachments & Список моделей вложений (связь с таблицей file) \\ \hline
Folder & Модель папки, в которой находится письмо \\ \hline
\end{xltabular}

\begin{xltabular}{\textwidth}{|p{6cm}|X|}
\caption{Методы класса MailMessageController\label{mailmessagecontroller:methods}} \\ \hline
\centrow Метод & \centrow Описание \\ \hline
\endfirsthead
\continuecaption{Продолжение таблицы \ref{mailmessagecontroller:methods}} 
\centrow Метод & \centrow Описание \\ \hline
\endhead
composeMail() & Открывает редактор письма с возможностью прикрепления файлов и выбора адресатов \\ \hline
sendMail() & Отправляет письмо через JMAP, проверяет лимиты вложений (10 МБ) \\ \hline
getMail() & Возвращает письмо по ID с проверкой прав доступа пользователя \\ \hline
deleteMail() & Помечает письмо как удалённое или перемещает в корзину \\ \hline
moveToFolder() & Изменяет папку письма (например, "Архив" → "Важные") \\ \hline
\end{xltabular}

\subsubsection{Спецификация класса MailFolderService}

Описание: Позволяет пользователю создавать, удалять и переименовывать папки для организации писем. Управляет системными папками ("Входящие", "Черновики"), обеспечивает синхронизацию с почтовым сервером и контроль квот на хранение.

Зависимости:
\begin{itemize}
    \item каждая папка принадлежит одному пользователю (связь с таблицей employee);
    \item содержит письма через связь с таблицей letter;
    \item интегрируется с MailMessageController для операций перемещения писем.
\end{itemize}

\begin{xltabular}{\textwidth}{|p{6cm}|X|}
\caption{Данные класса MailFolderService\label{mailfolderservice:data}} \\ \hline
\centrow Данные & \centrow Описание \\ \hline
\endfirsthead
\continuecaption{Продолжение таблицы \ref{mailfolderservice:data}} 
\centrow Данные & \centrow Описание \\ \hline
\endhead
Folder & Модель, представляющая папку \\ \hline
Name & Название папки (кириллица/латиница) \\ \hline
User & Владелец папки (связь с employee) \\ \hline
Type & Тип папки: "system" (входящие/отправленные) или "user" \\ \hline
\end{xltabular}

\begin{xltabular}{\textwidth}{|p{6cm}|X|}
\caption{Методы класса MailFolderService\label{mailfolderservice:methods}} \\ \hline
\centrow Метод & \centrow Описание \\ \hline
\endfirsthead
\continuecaption{Продолжение таблицы \ref{mailfolderservice:methods}} 
\centrow Метод & \centrow Описание \\ \hline
\endhead
createFolder() & Создаёт папку с валидацией имени (запрещены спецсимволы) \\ \hline
renameFolder() & Переименовывает папку, если она не системная \\ \hline
deleteFolder() & Удаляет папку после проверки на пустоту \\ \hline
listFolders() & Возвращает список папок с иерархией и статистикой \\ \hline
\end{xltabular}

\subsubsection{Спецификация класса ProjectManagerController}

Описание: Управляет созданием, редактированием и удалением проектов, настройкой целей, сроков и состава команды. Обеспечивает распределение задач между участниками, ведёт учёт этапов выполнения и позволяет просматривать общее состояние проекта.

Зависимости:
\begin{itemize}
    \item проект связан с пользователями через таблицу employee\_project;
    \item каждая задача проекта связывается через таблицу task;
    \item используются модули CalendarService и FileStorageService для назначения сроков и прикрепления файлов;
    \item интеграция с NotificationService для рассылки обновлений по задачам.
\end{itemize}

\begin{xltabular}{\textwidth}{|p{6cm}|X|}
\caption{Данные класса ProjectManagerController\label{projectmanager:data}} \\ \hline
\centrow Данные & \centrow Описание \\ \hline
\endfirsthead
\continuecaption{Продолжение таблицы \ref{projectmanager:data}}
\centrow Данные & \centrow Описание \\ \hline
\endhead
Project & Модель проекта с полями: название, описание, дедлайн \\ \hline
TeamMembers & Список сотрудников, назначенных на проект \\ \hline
Tasks & Список задач, связанных с проектом \\ \hline
Status & Текущий статус проекта: Active / On hold / Completed \\ \hline
Progress & Расчётный процент выполнения проекта \\ \hline
\end{xltabular}

\begin{xltabular}{\textwidth}{|p{6cm}|X|}
\caption{Методы класса ProjectManagerController\label{projectmanager:methods}} \\ \hline
\centrow Метод & \centrow Описание \\ \hline
\endfirsthead
\continuecaption{Продолжение таблицы \ref{projectmanager:methods}}
\centrow Метод & \centrow Описание \\ \hline
\endhead
createProject() & Создаёт новый проект с привязкой участников и сроков \\ \hline
editProject() & Редактирует основные параметры проекта \\ \hline
deleteProject() & Удаляет проект, если все задачи завершены \\ \hline
assignTeamMember() & Назначает сотрудника на проект \\ \hline
calculateProgress() & Вычисляет прогресс по завершённым задачам \\ \hline
\end{xltabular}

\subsubsection{Спецификация класса TaskTrackingService}

Описание: Отвечает за управление задачами в рамках проекта — создание, назначение, контроль выполнения, добавление комментариев и файлов. Обеспечивает гибкое управление статусами задач и контроль сроков выполнения.

Зависимости:
\begin{itemize}
    \item задачи связаны с проектами через поле project\_id;
    \item каждая задача может быть связана с несколькими исполнителями;
    \item интеграция с FileStorageService и CalendarController;
    \item подключается к NotificationService для напоминаний и уведомлений.
\end{itemize}

\begin{xltabular}{\textwidth}{|p{6cm}|X|}
\caption{Данные класса TaskTrackingService\label{tasktracking:data}} \\ \hline
\centrow Данные & \centrow Описание \\ \hline
\endfirsthead
\continuecaption{Продолжение таблицы \ref{tasktracking:data}}
\centrow Данные & \centrow Описание \\ \hline
\endhead
Task & Модель задачи с полями: описание, приоритет, статус \\ \hline
Assignee & Назначенный сотрудник или группа \\ \hline
Deadline & Крайний срок выполнения \\ \hline
Status & Этап выполнения: Todo / In Progress / Done \\ \hline
Comments & Список комментариев по задаче \\ \hline
\end{xltabular}

\begin{xltabular}{\textwidth}{|p{6cm}|X|}
\caption{Методы класса TaskTrackingService\label{tasktracking:methods}} \\ \hline
\centrow Метод & \centrow Описание \\ \hline
\endfirsthead
\continuecaption{Продолжение таблицы \ref{tasktracking:methods}}
\centrow Метод & \centrow Описание \\ \hline
\endhead
createTask() & Создание новой задачи в рамках проекта \\ \hline
updateStatus() & Обновление этапа выполнения задачи \\ \hline
addComment() & Добавление текстового комментария к задаче \\ \hline
assignEmployee() & Назначение исполнителя на задачу \\ \hline
attachFile() & Прикрепление файла к задаче \\ \hline
\end{xltabular}

\subsubsection{Спецификация класса CalendarController}

Описание: Управляет созданием, редактированием, отображением и удалением событий. Интегрируется с модулем видеоконференций для автоматического добавления ссылок на встречи. Поддерживает повторяющиеся события, приглашения участников и визуализацию в календарной сетке.

Зависимости:
\begin{itemize}
    \item события связаны с пользователями (создатель и участники) через таблицу employee;
    \item каждое событие может быть связано с видеовстречей через VKSController;
    \item взаимодействие с NotificationService для отправки приглашений и напоминаний;
    \item интеграция с ProjectManagerController для отображения проектных дедлайнов.
\end{itemize}

\begin{xltabular}{\textwidth}{|p{6cm}|X|}
\caption{Данные класса CalendarController\label{calendarcontroller:data}} \\ \hline
\centrow Данные & \centrow Описание \\ \hline
\endfirsthead
\continuecaption{Продолжение таблицы \ref{calendarcontroller:data}}
\centrow Данные & \centrow Описание \\ \hline
\endhead
Event & Модель события с полями: тема, дата, время, описание \\ \hline
Creator & Сотрудник, создавший событие \\ \hline
Participants & Список участников события \\ \hline
RecurrenceRule & Правила повторения события (ежедневно/еженедельно) \\ \hline
ConferenceLink & Ссылка на видеоконференцию (если доступно) \\ \hline
\end{xltabular}

\begin{xltabular}{\textwidth}{|p{6cm}|X|}
\caption{Методы класса CalendarController\label{calendarcontroller:methods}} \\ \hline
\centrow Метод & \centrow Описание \\ \hline
\endfirsthead
\continuecaption{Продолжение таблицы \ref{calendarcontroller:methods}}
\centrow Метод & \centrow Описание \\ \hline
\endhead
createEvent() & Создаёт событие с указанием даты, времени и участников \\ \hline
updateEvent() & Редактирует параметры события и уведомляет участников \\ \hline
deleteEvent() & Удаляет событие и очищает связанные встречи \\ \hline
syncWithVKS() & Интегрирует событие с видеовстречей \\ \hline
getEventList() & Возвращает список событий для пользователя \\ \hline
\end{xltabular}



\subsubsection{Спецификация класса ReminderService}

Описание: Отвечает за отправку уведомлений пользователям о предстоящих событиях. Поддерживает как email-напоминания, так и push-уведомления. Может быть связан с другими модулями (проекты, встречи) для централизованного оповещения.

Зависимости:
\begin{itemize}
    \item события загружаются из CalendarController;
    \item использует NotificationService для доставки уведомлений;
    \item доступ к пользователям через таблицу employee;
    \item может быть переиспользован для ProjectManagerController.
\end{itemize}

\begin{xltabular}{\textwidth}{|p{6cm}|X|}
\caption{Данные класса ReminderService\label{reminderservice:data}} \\ \hline
\centrow Данные & \centrow Описание \\ \hline
\endfirsthead
\continuecaption{Продолжение таблицы \ref{reminderservice:data}}
\centrow Данные & \centrow Описание \\ \hline
\endhead
Reminder & Модель напоминания (тип, время, канал) \\ \hline
User & Сотрудник, для которого установлено напоминание \\ \hline
EventTime & Время начала события, к которому привязано напоминание \\ \hline
DeliveryChannel & Канал отправки: email / push / system \\ \hline
IsAcknowledged & Статус подтверждения уведомления пользователем \\ \hline
\end{xltabular}

\begin{xltabular}{\textwidth}{|p{6cm}|X|}
\caption{Методы класса ReminderService\label{reminderservice:methods}} \\ \hline
\centrow Метод & \centrow Описание \\ \hline
\endfirsthead
\continuecaption{Продолжение таблицы \ref{reminderservice:methods}}
\centrow Метод & \centrow Описание \\ \hline
\endhead
setReminder() & Устанавливает напоминание по дате и времени \\ \hline
sendNotifications() & Рассылает уведомления участникам события \\ \hline
checkUpcomingEvents() & Проверяет и активирует напоминания за 15 мин. до события \\ \hline
acknowledgeReminder() & Отмечает напоминание как просмотренное \\ \hline
deleteReminder() & Удаляет напоминание по ID или событию \\ \hline
\end{xltabular}

\subsubsection{Спецификация класса FileManager}

Описание: Отвечает за операции с файлами и папками: загрузку, скачивание, переименование, удаление, создание структуры директорий. Контролирует права доступа к объектам, работает с превью изображений, отображает использование дискового пространства. Используется в модулях «Проекты», «Почта», «Разговоры».

Зависимости:
\begin{itemize}
    \item каждый файл связан с пользователем (через employee) и папкой (через folder);
    \item работает с FileStorageService для хранения и шифрования данных;
    \item использует PermissionService для контроля доступа;
    \item интеграция с TaskTrackingService и MailMessageController для вложений.
\end{itemize}

\begin{xltabular}{\textwidth}{|p{6cm}|X|}
\caption{Данные класса FileManager\label{filemanager:data}} \\ \hline
\centrow Данные & \centrow Описание \\ \hline
\endfirsthead
\continuecaption{Продолжение таблицы \ref{filemanager:data}}
\centrow Данные & \centrow Описание \\ \hline
\endhead
File & Модель файла (имя, размер, расширение, hash) \\ \hline
Owner & Владелец файла (ID пользователя) \\ \hline
Folder & Папка, в которой находится файл \\ \hline
AccessRights & Уровень доступа: read/write/delete \\ \hline
Preview & Ссылка на превью (если доступно) \\ \hline
CreatedAt & Дата загрузки файла \\ \hline
\end{xltabular}

\begin{xltabular}{\textwidth}{|p{6cm}|X|}
\caption{Методы класса FileManager\label{filemanager:methods}} \\ \hline
\centrow Метод & \centrow Описание \\ \hline
\endfirsthead
\continuecaption{Продолжение таблицы \ref{filemanager:methods}}
\centrow Метод & \centrow Описание \\ \hline
\endhead
uploadFile() & Загружает файл, проверяя тип и размер \\ \hline
downloadFile() & Скачивает файл по ID, проверяя права доступа \\ \hline
deleteFile() & Удаляет файл с сервера и БД \\ \hline
createFolder() & Создаёт новую директорию для хранения файлов \\ \hline
moveFile() & Перемещает файл между папками \\ \hline
\end{xltabular}



\subsubsection{Спецификация класса FileStorageService}

Описание: Обеспечивает безопасное хранение файлов, отвечает за шифрование, резервное копирование, проверку целостности и антивирусную проверку. Является низкоуровневым слоем между FileManager и физическим хранилищем.

Зависимости:
\begin{itemize}
    \item получает запросы от FileManager для операций над файлами;
    \item использует внешние API для шифрования (например, AES256);
    \item подключает антивирусный движок ClamAV (или аналог);
    \item взаимодействует с BackupService (резервные копии).
\end{itemize}

\begin{xltabular}{\textwidth}{|p{6cm}|X|}
\caption{Данные класса FileStorageService\label{filestorageservice:data}} \\ \hline
\centrow Данные & \centrow Описание \\ \hline
\endfirsthead
\continuecaption{Продолжение таблицы \ref{filestorageservice:data}}
\centrow Данные & \centrow Описание \\ \hline
\endhead
StoragePath & Путь к директории хранения файлов \\ \hline
EncryptionKey & Ключ для шифрования содержимого файла \\ \hline
Checksum & Контрольная сумма файла для проверки целостности \\ \hline
ScanStatus & Статус антивирусной проверки \\ \hline
BackupTimestamp & Время последней резервной копии \\ \hline
\end{xltabular}

\begin{xltabular}{\textwidth}{|p{6cm}|X|}
\caption{Методы класса FileStorageService\label{filestorageservice:methods}} \\ \hline
\centrow Метод & \centrow Описание \\ \hline
\endfirsthead
\continuecaption{Продолжение таблицы \ref{filestorageservice:methods}}
\centrow Метод & \centrow Описание \\ \hline
\endhead
encryptFile() & Шифрует файл перед сохранением \\ \hline
decryptFile() & Расшифровывает файл при загрузке \\ \hline
scanForViruses() & Проверяет файл на вирусы перед сохранением \\ \hline
createBackup() & Сохраняет резервную копию файла в отдельный storage \\ \hline
verifyChecksum() & Проверяет контрольную сумму при загрузке \\ \hline
\end{xltabular}

\subsubsection{Спецификация класса VKSController}

Описание: Отвечает за организацию видеоконференций между пользователями платформы. Управляет созданием комнат, генерацией ссылок на подключение, настройкой доступа, записью и завершением встреч. Обеспечивает привязку к событиям календаря и проектным задачам.

Зависимости:
\begin{itemize}
    \item встречи связаны с пользователями через таблицу employee;
    \item каждая конференция может быть связана с событием в календаре;
    \item взаимодействие с MediaService для передачи видеопотоков и записи;
    \item интеграция с NotificationService для рассылки приглашений.
\end{itemize}

\begin{xltabular}{\textwidth}{|p{6cm}|X|}
\caption{Данные класса VKSController\label{vkscontroller:data}} \\ \hline
\centrow Данные & \centrow Описание \\ \hline
\endfirsthead
\continuecaption{Продолжение таблицы \ref{vkscontroller:data}}
\centrow Данные & \centrow Описание \\ \hline
\endhead
Conference & Модель конференции (ID, тема, дата, продолжительность) \\ \hline
Host & Ведущий (создатель конференции) \\ \hline
Participants & Приглашённые пользователи (список employee ID) \\ \hline
AccessCode & Код доступа (если требуется) \\ \hline
RecordingEnabled & Флаг записи встречи \\ \hline
LinkedEvent & Событие календаря, связанное с конференцией \\ \hline
\end{xltabular}

\begin{xltabular}{\textwidth}{|p{6cm}|X|}
\caption{Методы класса VKSController\label{vkscontroller:methods}} \\ \hline
\centrow Метод & \centrow Описание \\ \hline
\endfirsthead
\continuecaption{Продолжение таблицы \ref{vkscontroller:methods}}
\centrow Метод & \centrow Описание \\ \hline
\endhead
createConference() & Создание конференции с параметрами доступа \\ \hline
inviteParticipants() & Отправка приглашений участникам \\ \hline
startConference() & Активация видеопотока и подключение пользователей \\ \hline
endConference() & Завершение встречи и выгрузка записи \\ \hline
generateJoinLink() & Генерация ссылки с токеном доступа \\ \hline
\end{xltabular}



\subsubsection{Спецификация класса MediaService}

Описание: Обеспечивает обработку аудио- и видеопотоков во время конференции. Управляет началом/окончанием трансляции, записью, сжатием и сохранением медиаданных. Также отвечает за качество соединения и адаптацию потока под пропускную способность.

Зависимости:
\begin{itemize}
    \item работает напрямую с VKSController через WebRTC API;
    \item хранит записи в FileStorageService;
    \item использует адаптивные алгоритмы качества (например, VP9, Opus);
    \item взаимодействует с нагрузочным балансировщиком при высоком трафике.
\end{itemize}

\begin{xltabular}{\textwidth}{|p{6cm}|X|}
\caption{Данные класса MediaService\label{mediaservice:data}} \\ \hline
\centrow Данные & \centrow Описание \\ \hline
\endfirsthead
\continuecaption{Продолжение таблицы \ref{mediaservice:data}}
\centrow Данные & \centrow Описание \\ \hline
\endhead
Stream & Видеопоток, получаемый по WebRTC \\ \hline
Bitrate & Текущая скорость передачи (kbps) \\ \hline
AudioCodec & Кодек звука (например, Opus) \\ \hline
VideoCodec & Кодек видео (например, VP9) \\ \hline
RecordingPath & Путь для сохранения записи встречи \\ \hline
Latency & Задержка при передаче (мс) \\ \hline
\end{xltabular}

\begin{xltabular}{\textwidth}{|p{6cm}|X|}
\caption{Методы класса MediaService\label{mediaservice:methods}} \\ \hline
\centrow Метод & \centrow Описание \\ \hline
\endfirsthead
\continuecaption{Продолжение таблицы \ref{mediaservice:methods}}
\centrow Метод & \centrow Описание \\ \hline
\endhead
setupStream() & Настройка видеопотока на стороне клиента \\ \hline
startRecording() & Начало записи медиапотока \\ \hline
stopRecording() & Завершение записи и сохранение файла \\ \hline
adjustBitrate() & Автоматическая настройка качества в зависимости от сети \\ \hline
monitorLatency() & Отслеживание задержки в реальном времени \\ \hline
\end{xltabular}

\subsubsection{Спецификация класса ChatController}

Описание: Управляет созданием чатов, добавлением/удалением участников, организацией каналов, групповых и приватных бесед. Поддерживает отображение истории, закрепление сообщений, а также интеграцию с проектами и контактами.

Зависимости:
\begin{itemize}
    \item чаты связаны с пользователями через таблицу employee;
    \item каждое сообщение привязано к чату через message.chat\_id;
    \item использует ContactController для отображения информации об участниках;
    \item взаимодействует с NotificationService и MessageService.
\end{itemize}

\begin{xltabular}{\textwidth}{|p{6cm}|X|}
\caption{Данные класса ChatController\label{chatcontroller:data}} \\ \hline
\centrow Данные & \centrow Описание \\ \hline
\endfirsthead
\continuecaption{Продолжение таблицы \ref{chatcontroller:data}}
\centrow Данные & \centrow Описание \\ \hline
\endhead
Chat & Модель чата (ID, название, тип) \\ \hline
Members & Список участников чата \\ \hline
IsGroup & Признак группового чата \\ \hline
PinnedMessages & Закреплённые сообщения (список ID) \\ \hline
CreatedAt & Дата создания чата \\ \hline
\end{xltabular}

\begin{xltabular}{\textwidth}{|p{6cm}|X|}
\caption{Методы класса ChatController\label{chatcontroller:methods}} \\ \hline
\centrow Метод & \centrow Описание \\ \hline
\endfirsthead
\continuecaption{Продолжение таблицы \ref{chatcontroller:methods}}
\centrow Метод & \centrow Описание \\ \hline
\endhead
createChat() & Создание нового чата (группового или приватного) \\ \hline
addMember() & Добавление участника в чат \\ \hline
removeMember() & Удаление участника из чата \\ \hline
pinMessage() & Закрепление сообщения в верхней части чата \\ \hline
getChatHistory() & Получение истории сообщений по chat\_id \\ \hline
\end{xltabular}



\subsubsection{Спецификация класса MessageService}

Описание: Обрабатывает отправку, получение и хранение сообщений. Обеспечивает доставку push-уведомлений, шифрование, поддержку вложений, реакций и статусов доставки (прочитано/доставлено).

Зависимости:
\begin{itemize}
    \item каждое сообщение связано с пользователем (author\_id) и чатом (chat\_id);
    \item вложения реализуются через FileManager;
    \item интеграция с NotificationService и EncryptionModule;
    \item статусы доставки хранятся в таблице message\_status.
\end{itemize}

\begin{xltabular}{\textwidth}{|p{6cm}|X|}
\caption{Данные класса MessageService\label{messageservice:data}} \\ \hline
\centrow Данные & \centrow Описание \\ \hline
\endfirsthead
\continuecaption{Продолжение таблицы \ref{messageservice:data}}
\centrow Данные & \centrow Описание \\ \hline
\endhead
Message & Текст сообщения, время отправки, тип (text/image/file) \\ \hline
Author & Отправитель сообщения (employee ID) \\ \hline
Chat & Чат, к которому относится сообщение \\ \hline
DeliveryStatus & Статус доставки (отправлено, доставлено, прочитано) \\ \hline
Attachments & Вложенные файлы (связь с FileManager) \\ \hline
\end{xltabular}

\begin{xltabular}{\textwidth}{|p{6cm}|X|}
\caption{Методы класса MessageService\label{messageservice:methods}} \\ \hline
\centrow Метод & \centrow Описание \\ \hline
\endfirsthead
\continuecaption{Продолжение таблицы \ref{messageservice:methods}}
\centrow Метод & \centrow Описание \\ \hline
\endhead
sendMessage() & Отправка сообщения в чат \\ \hline
receiveMessage() & Получение входящего сообщения по WebSocket \\ \hline
deleteMessage() & Удаление сообщения (soft delete) \\ \hline
markAsRead() & Обновление статуса доставки до "прочитано" \\ \hline
encryptMessage() & Шифрование тела сообщения перед отправкой \\ \hline
\end{xltabular}

\subsubsection{Спецификация класса ContactController}

Описание: Управляет добавлением, редактированием, удалением и группировкой контактов. Обеспечивает интеграцию с почтой, чатами и системой уведомлений. Поддерживает метки, фильтрацию, сортировку и массовый импорт из внешних источников.

Зависимости:
\begin{itemize}
    \item каждый контакт принадлежит пользователю (employee);
    \item группы контактов реализованы через таблицу contact\_group;
    \item интеграция с MailMessageController для автозаполнения email;
    \item доступ к контактам ограничен ролями и политиками приватности.
\end{itemize}

\begin{xltabular}{\textwidth}{|p{6cm}|X|}
\caption{Данные класса ContactController\label{contactcontroller:data}} \\ \hline
\centrow Данные & \centrow Описание \\ \hline
\endfirsthead
\continuecaption{Продолжение таблицы \ref{contactcontroller:data}}
\centrow Данные & \centrow Описание \\ \hline
\endhead
Contact & Модель контакта (имя, email, телефон, компания) \\ \hline
Group & Группа контактов, к которой относится запись \\ \hline
Label & Метки для фильтрации (JSON-массив) \\ \hline
IsFavorite & Признак избранного контакта \\ \hline
CreatedAt & Дата добавления контакта \\ \hline
\end{xltabular}

\begin{xltabular}{\textwidth}{|p{6cm}|X|}
\caption{Методы класса ContactController\label{contactcontroller:methods}} \\ \hline
\centrow Метод & \centrow Описание \\ \hline
\endfirsthead
\continuecaption{Продолжение таблицы \ref{contactcontroller:methods}}
\centrow Метод & \centrow Описание \\ \hline
\endhead
addContact() & Добавление нового контакта в справочник \\ \hline
editContact() & Редактирование данных существующего контакта \\ \hline
deleteContact() & Удаление контакта с проверкой связей \\ \hline
groupContacts() & Группировка по категориям \\ \hline
getContactList() & Получение списка контактов пользователя \\ \hline
\end{xltabular}



\subsubsection{Спецификация класса ContactSyncService}

Описание: Обеспечивает синхронизацию контактов между локальной базой и внешними источниками: почтовыми клиентами, CRM-системами, корпоративными LDAP-серверами. Поддерживает очистку дубликатов и регулярные автообновления.

Зависимости:
\begin{itemize}
    \item контакты синхронизируются с внешними источниками (Google, Outlook, LDAP);
    \item взаимодействует с ContactController и NotificationService;
    \item использует алгоритмы сравнения строк (например, Levenshtein);
    \item поддерживает экспорт в формате vCard и CSV.
\end{itemize}

\begin{xltabular}{\textwidth}{|p{6cm}|X|}
\caption{Данные класса ContactSyncService\label{contactsyncservice:data}} \\ \hline
\centrow Данные & \centrow Описание \\ \hline
\endfirsthead
\continuecaption{Продолжение таблицы \ref{contactsyncservice:data}}
\centrow Данные & \centrow Описание \\ \hline
\endhead
ExternalSource & Информация об источнике данных (тип, URL) \\ \hline
LastSync & Время последней синхронизации \\ \hline
ConflictLog & Список конфликтов при объединении записей \\ \hline
DuplicateCount & Количество найденных дубликатов \\ \hline
Format & Формат импорта/экспорта (CSV/vCard/LDAP) \\ \hline
\end{xltabular}

\begin{xltabular}{\textwidth}{|p{6cm}|X|}
\caption{Методы класса ContactSyncService\label{contactsyncservice:methods}} \\ \hline
\centrow Метод & \centrow Описание \\ \hline
\endfirsthead
\continuecaption{Продолжение таблицы \ref{contactsyncservice:methods}}
\centrow Метод & \centrow Описание \\ \hline
\endhead
importContacts() & Импортирует контакты из указанного источника \\ \hline
exportContacts() & Экспортирует контакты в выбранном формате \\ \hline
deduplicate() & Находит и объединяет дублирующиеся записи \\ \hline
autoSync() & Планирует регулярную синхронизацию \\ \hline
resolveConflicts() & Обрабатывает коллизии при совпадении данных \\ \hline
\end{xltabular}

\subsubsection{Спецификация класса SettingsController}

Описание: Управляет пользовательскими настройками: профилем, паролем, темой оформления, уведомлениями и языком интерфейса. Является точкой входа для фронтенда при инициализации настроек в сессии пользователя. Также осуществляет базовую валидацию и сохранение изменений.

Зависимости:
\begin{itemize}
    \item данные сохраняются в PreferencesService;
    \item используется AuthService для смены пароля и получения информации о пользователе;
    \item взаимодействует с NotificationService для настройки уведомлений;
    \item интеграция с интерфейсной темой (через ThemeProvider).
\end{itemize}

\begin{xltabular}{\textwidth}{|p{6cm}|X|}
\caption{Данные класса SettingsController\label{settingscontroller:data}} \\ \hline
\centrow Данные & \centrow Описание \\ \hline
\endfirsthead
\continuecaption{Продолжение таблицы \ref{settingscontroller:data}}
\centrow Данные & \centrow Описание \\ \hline
\endhead
UserProfile & Информация о пользователе (ФИО, email, аватар) \\ \hline
Theme & Выбранная тема (light/dark) \\ \hline
NotificationSettings & Типы включённых уведомлений \\ \hline
Language & Текущий язык интерфейса \\ \hline
TwoFactorEnabled & Признак включённой двухфакторной аутентификации \\ \hline
\end{xltabular}

\begin{xltabular}{\textwidth}{|p{6cm}|X|}
\caption{Методы класса SettingsController\label{settingscontroller:methods}} \\ \hline
\centrow Метод & \centrow Описание \\ \hline
\endfirsthead
\continuecaption{Продолжение таблицы \ref{settingscontroller:methods}}
\centrow Метод & \centrow Описание \\ \hline
\endhead
updateProfile() & Обновление информации пользователя \\ \hline
changePassword() & Изменение пароля после проверки старого \\ \hline
setTheme() & Установка светлой или тёмной темы \\ \hline
configureNotifications() & Настройка уведомлений (email/push) \\ \hline
toggle2FA() & Включение или отключение двухфакторной авторизации \\ \hline
\end{xltabular}

\subsubsection{Спецификация класса PreferencesService}

Описание: Обеспечивает сохранение и загрузку настроек пользователя. Хранит значения в привязанной таблице настроек в БД или в localStorage при offline-режиме. Отвечает за дефолтные параметры и восстановление настроек при сбросе.

Зависимости:
\begin{itemize}
    \item используется SettingsController как интерфейс управления;
    \item взаимодействует с AuthService для получения user\_id;
    \item использует таблицу preferences (user\_id, key, value);
    \item подключается к DefaultSettingsProvider при первом входе.
\end{itemize}

\begin{xltabular}{\textwidth}{|p{6cm}|X|}
\caption{Данные класса PreferencesService\label{preferencesservice:data}} \\ \hline
\centrow Данные & \centrow Описание \\ \hline
\endfirsthead
\continuecaption{Продолжение таблицы \ref{preferencesservice:data}}
\centrow Данные & \centrow Описание \\ \hline
\endhead
UserId & Идентификатор пользователя (внешний ключ) \\ \hline
Key & Название параметра (например, theme) \\ \hline
Value & Значение параметра (например, dark) \\ \hline
UpdatedAt & Время последнего изменения \\ \hline
IsDefault & Флаг, указывающий на дефолтность значения \\ \hline
\end{xltabular}

\begin{xltabular}{\textwidth}{|p{6cm}|X|}
\caption{Методы класса PreferencesService\label{preferencesservice:methods}} \\ \hline
\centrow Метод & \centrow Описание \\ \hline
\endfirsthead
\continuecaption{Продолжение таблицы \ref{preferencesservice:methods}}
\centrow Метод & \centrow Описание \\ \hline
\endhead
saveSetting() & Сохраняет или обновляет параметр настройки \\ \hline
loadSettings() & Загружает все настройки пользователя по user\_id \\ \hline
resetToDefault() & Сбрасывает настройки к значениям по умолчанию \\ \hline
getSetting() & Получает конкретное значение по ключу \\ \hline
mergeSettings() & Объединяет локальные и серверные настройки \\ \hline
\end{xltabular}

\subsubsection{Спецификация класса DashboardController}

Описание: Отвечает за отображение пользовательской панели, управление расположением и конфигурацией виджетов. Позволяет добавлять, удалять и перетаскивать виджеты, обновлять данные в реальном времени. Также обеспечивает сохранение пользовательского расположения и адаптацию под устройства.

Зависимости:
\begin{itemize}
    \item виджеты связаны с пользователями через таблицу dashboard\_layout;
    \item получает данные с других сервисов: TaskTrackingService, FileManager и др.;
    \item используется WidgetService для генерации и обновления содержимого;
    \item взаимодействует с PreferencesService для сохранения расположения.
\end{itemize}

\begin{xltabular}{\textwidth}{|p{6cm}|X|}
\caption{Данные класса DashboardController\label{dashboardcontroller:data}} \\ \hline
\centrow Данные & \centrow Описание \\ \hline
\endfirsthead
\continuecaption{Продолжение таблицы \ref{dashboardcontroller:data}}
\centrow Данные & \centrow Описание \\ \hline
\endhead
UserLayout & Массив с позицией и конфигурацией виджетов \\ \hline
AvailableWidgets & Список всех доступных виджетов платформы \\ \hline
ActiveWidgets & Список активных виджетов для пользователя \\ \hline
RefreshInterval & Интервал обновления данных на панели \\ \hline
GridSettings & Настройки сетки (столбцы, отступы, адаптивность) \\ \hline
\end{xltabular}

\begin{xltabular}{\textwidth}{|p{6cm}|X|}
\caption{Методы класса DashboardController\label{dashboardcontroller:methods}} \\ \hline
\centrow Метод & \centrow Описание \\ \hline
\endfirsthead
\continuecaption{Продолжение таблицы \ref{dashboardcontroller:methods}}
\centrow Метод & \centrow Описание \\ \hline
\endhead
addWidget() & Добавляет выбранный виджет на панель пользователя \\ \hline
removeWidget() & Удаляет виджет с панели и очищает его конфигурацию \\ \hline
updateLayout() & Обновляет позиции виджетов после перемещения \\ \hline
loadDashboard() & Загружает текущую конфигурацию панели \\ \hline
refreshData() & Принудительное обновление содержимого всех виджетов \\ \hline
\end{xltabular}

\subsubsection{Спецификация класса WidgetService}

Описание: Обрабатывает генерацию, конфигурацию и рендеринг виджетов на панели. Отвечает за запрос и кэширование данных, отображение графиков, таблиц, уведомлений и прочих блоков. Позволяет адаптировать содержимое под роль пользователя.

Зависимости:
\begin{itemize}
    \item данные поступают от сторонних сервисов: TaskTrackingService, FileManager и т.д.;
    \item виджеты отображаются через DashboardController;
    \item использует PreferencesService для хранения индивидуальных настроек;
    \item взаимодействует с ChartRenderer, TableComponent, HTMLWidget.
\end{itemize}

\begin{xltabular}{\textwidth}{|p{6cm}|X|}
\caption{Данные класса WidgetService\label{widgetservice:data}} \\ \hline
\centrow Данные & \centrow Описание \\ \hline
\endfirsthead
\continuecaption{Продолжение таблицы \ref{widgetservice:data}}
\centrow Данные & \centrow Описание \\ \hline
\endhead
Widget & Модель виджета (тип, параметры, источник данных) \\ \hline
Config & Конфигурация отображения (цвет, размер, тип графика) \\ \hline
DataCache & Кэш последних полученных данных \\ \hline
Permissions & Уровни доступа к данным виджета \\ \hline
UpdateFrequency & Частота обновления содержимого \\ \hline
\end{xltabular}

\begin{xltabular}{\textwidth}{|p{6cm}|X|}
\caption{Методы класса WidgetService\label{widgetservice:methods}} \\ \hline
\centrow Метод & \centrow Описание \\ \hline
\endfirsthead
\continuecaption{Продолжение таблицы \ref{widgetservice:methods}}
\centrow Метод & \centrow Описание \\ \hline
\endhead
createWidget() & Инициализирует новый виджет по типу \\ \hline
updateWidget() & Обновляет конфигурацию и источник данных \\ \hline
fetchData() & Получает актуальные данные от соответствующего сервиса \\ \hline
renderWidget() & Отображает содержимое в нужном формате (таблица, график и т.д.) \\ \hline
cacheData() & Сохраняет последние данные в кэше \\ \hline
\end{xltabular}

\subsection{Тестирование программной системы}

\subsubsection{Модульное тестирование классов и методов программы}

Модульное тестирование в рамках разработки веб-платформы проводится с целью проверки корректной работы каждого отдельного класса и метода системы. Оно позволяет локализовать и устранить ошибки ещё на ранней стадии, обеспечивая надёжность и предсказуемое поведение компонентов при различных сценариях использования.

Описание процесса модульного тестирования:

Модульные тесты разрабатываются для всех основных бизнес-классов, отвечающих за управление почтой, контактами, задачами, видеоконференциями, файлами, настройками и панелью управления. Для тестирования используется автоматизированный фреймворк (например, Jest для компонентов на TypeScript или pytest для Python-частей), позволяющий запускать тесты изолированно и проверять возвращаемые значения, обработку ошибок, работу с внешними сервисами и корректность изменений состояния.

Цели модульного тестирования:
\begin{enumerate}
    \item Изоляция — каждый тест проверяет только один метод или функциональный блок.
    \item Автоматизация — тесты можно запускать непрерывно (CI), без вмешательства пользователя.
    \item Покрытие — тестируются все ключевые методы, включая граничные случаи и обработку исключений.
\end{enumerate}

\newpage
Тестирование MailMessageController:

\begin{figure}[H]
\begin{lstlisting}[language=JavaScript]
import { MailMessageController } from '../controllers/MailMessageController';
import { mockJMAPService } from '../mocks/jmapService';
import { mockStorage } from '../mocks/storage';

describe('MailMessageController', () => {
  let controller;

  beforeEach(() => {
    controller = new MailMessageController(mockJMAPService, mockStorage);
  });

  it('успешно отправляет письмо с корректными данными', async () => {
    const result = await controller.sendMail({
      subject: 'Тестовое письмо',
      body: 'Содержимое письма',
      recipients: ['test@example.com'],
      attachments: [],
    });

    expect(result.status).toBe('success');
    expect(result.messageId).toBeDefined();
  });

  it('выдаёт ошибку при отсутствии получателя', async () => {
    await expect(
      controller.sendMail({
        subject: 'Без получателя',
        body: 'Текст',
        recipients: [],
        attachments: [],
      })
    ).rejects.toThrow('Поле "recipients" обязательно');
  });
});
\end{lstlisting}
\captionof{figure}{Тесты TaskScheduler}
\end{figure}

\newpage
Тестирование MailMessageController:

\begin{figure}[H]
\begin{lstlisting}[language=Python]
import pytest
from myapp.controllers import MailMessageController
from unittest.mock import Mock

@pytest.fixture
def mail_controller():
    mock_jmap_service = Mock()
    mock_storage = Mock()
    return MailMessageController(mock_jmap_service, mock_storage)

def test_send_mail_success(mail_controller):
    mail_controller.jmap_service.send.return_value = "msg123"
    result = mail_controller.send_mail(
        subject="Тестовое письмо",
        body="Содержимое письма",
        recipients=["test@example.com"],
        attachments=[]
    )
    assert result['status'] == 'success'
    assert 'message_id' in result

def test_send_mail_no_recipients(mail_controller):
    with pytest.raises(ValueError) as exc_info:
        mail_controller.send_mail(
            subject="Без получателя",
            body="Текст",
            recipients=[],
            attachments=[]
        )
    assert "recipients" in str(exc_info.value)
\end{lstlisting}
\captionof{figure}{Тесты MailMessageController}
\end{figure}


\newpage
Тестирование ContactManager:

\begin{figure}[H]
\begin{lstlisting}[language=JavaScript]
import { ContactManager } from '../controllers/ContactManager';

describe('ContactManager', () => {
  let manager;

  beforeEach(() => {
    manager = new ContactManager();
  });

  it('добавляет новый контакт', () => {
    const contact = { id: '1', name: 'Иван' };
    manager.addContact(contact);
    expect(manager.getContactById('1')).toEqual(contact);
  });

  it('удаляет контакт', () => {
    manager.addContact({ id: '2', name: 'Мария' });
    manager.removeContact('2');
    expect(manager.getContactById('2')).toBeUndefined();
  });
});
\end{lstlisting}
\captionof{figure}{Тесты ContactManager}
\end{figure}

\newpage
Тестирование TaskScheduler:

\begin{figure}[H]
\begin{lstlisting}[language=JavaScript]
import { TaskScheduler } from '../controllers/TaskScheduler';

describe('TaskScheduler', () => {
  let scheduler;

  beforeEach(() => {
    scheduler = new TaskScheduler();
    jest.useFakeTimers();
  });

  afterEach(() => {
    jest.useRealTimers();
  });

  it('запускает задачу по расписанию', () => {
    const mockTask = jest.fn();
    scheduler.schedule('task1', mockTask, 1000);
    jest.advanceTimersByTime(1000);
    expect(mockTask).toHaveBeenCalled();
  });

  it('отменяет задачу', () => {
    const mockTask = jest.fn();
    scheduler.schedule('task2', mockTask, 1000);
    scheduler.cancel('task2');
    jest.advanceTimersByTime(1000);
    expect(mockTask).not.toHaveBeenCalled();
  });
});
\end{lstlisting}
\captionof{figure}{Тесты TaskScheduler}
\end{figure}

\newpage
Тестирование VideoConferenceManager:

\begin{figure}[H]
\begin{lstlisting}[language=Python]
import pytest
from myapp.controllers import VideoConferenceManager

@pytest.fixture
def vc_manager():
    return VideoConferenceManager()

def test_start_meeting(vc_manager):
    meeting_id = vc_manager.start_meeting(['user1', 'user2'])
    assert isinstance(meeting_id, str)
    assert vc_manager.is_active(meeting_id)

def test_end_meeting(vc_manager):
    meeting_id = vc_manager.start_meeting(['user1'])
    vc_manager.end_meeting(meeting_id)
    assert not vc_manager.is_active(meeting_id)
\end{lstlisting}
\captionof{figure}{Тесты VideoConferenceManager}
\end{figure}

\newpage
Тестирование FileStorageController:

\begin{figure}[H]
\begin{lstlisting}[language=Python]
import pytest
from myapp.controllers import FileStorageController

@pytest.fixture
def storage_controller():
    return FileStorageController()

def test_save_file(storage_controller):
    file_id = storage_controller.save_file(b'content', 'file.txt')
    assert file_id is not None

def test_delete_file(storage_controller):
    file_id = storage_controller.save_file(b'content', 'file.txt')
    storage_controller.delete_file(file_id)
    assert not storage_controller.exists(file_id)
\end{lstlisting}
\captionof{figure}{Тесты FileStorageController}
\end{figure}

Результаты тестирования

Все представленные тесты успешно проверяют поведение ключевых компонентов приложения. Использование Jest и Pytest обеспечивает:
\begin{itemize}
  \item контроль корректности бизнес-логики;
  \item отлавливание ошибок при неверных входных данных;
  \item поддержку модульного и интеграционного тестирования.
\end{itemize}